\section{Introduction}

Learning-to-defer (L2D) studies decision systems that \emph{route} each query to one of several experts and incur expert-dependent
\emph{consultation costs} \citep{madras2018predict,mozannar2021consistent,Narasimhan,mao2023twostage, montreuil2025ask}.
Most L2D work is offline: a router is trained on a fixed dataset (often under i.i.d.\ assumptions) using supervision that is typically unavailable online,
such as access to all experts' predictions or losses for the same input.


In sequential settings, decisions and observations are interleaved over time, making the offline assumptions above unrealistic.
Such regimes arise, for instance, in streaming forecasting or triage pipelines where one must route each incoming instance to a single expert under limited and time-varying capacity.
At round $t$, the router observes a context $\mathbf{x}_t$ and the currently available expert set $\mathcal{E}_t$, selects $I_t \in \mathcal{E}_t$, and then observes the outcome $\vy_t$ along with the queried prediction $\vhaty_{t,I_t}$.
Feedback is \emph{partial}: predictions from $j\neq I_t$ are never revealed.
The process is typically \emph{non-i.i.d.} and often non-stationary \citep{hamilton2020time, sezer2020financial}, so expert quality and cross-expert dependencies may drift or switch across regimes.
In addition, $\mathcal{E}_t$ itself can evolve as experts arrive, depart, or become temporarily unavailable, and querying can consume scarce resources subject to operational constraints.
Altogether, partial feedback, distribution shift, and a time-varying expert pool render standard offline L2D formulations inadequate and motivate online methods that explicitly track uncertainty over time.



To address these challenges, we develop a probabilistic routing framework for non-stationary time series
under partial feedback and a dynamic expert pool.
We model expert residuals with a switching linear-Gaussian state-space
\citep{ghahramani2000variational, linderman2016recurrent, hu2024modeling} model
that couples a shared global factor with
expert-specific idiosyncratic states and a discrete regime process, enabling time-varying cross-expert dependence.
Faithful to practical settings, we support adding or removing experts without affecting the maintained marginals of retained experts.
We also propose an Information-Directed Sampling (IDS)-inspired exploration rule \citep{russo2014learning} that trades off predicted cost against information gained about the latent regime and shared factor.

\textbf{Contributions.}
Our main contributions are:
\begin{itemize}[topsep=2pt,itemsep=1pt,parsep=0pt,leftmargin=1.2em]
    \item We formalize \emph{online} expert routing for non-stationary time series with bandit feedback and time-varying expert availability (Section~\ref{sec:problem_formulation}).
    \item We propose \textbf{L2D-SLDS}, a factorized switching linear-Gaussian state-space model for expert residuals with context-dependent regime transitions, a shared global factor enabling cross-expert information transfer, and per-expert idiosyncratic dynamics (Section~\ref{sec:generative_model}).
    \item We derive an efficient Interacting Multiple Model (IMM)-style \citep{johnston2002improvement} filtering recursion under censoring together with dynamic registry management for expert entry and pruning; we show pruning leaves retained marginals unchanged (Sections~\ref{sec:generative_model} and~\ref{sec:registry}; Proposition~\ref{prop:invariance}).
    \item We introduce an IDS-inspired routing rule \citep{russo2014learning} that trades off predicted cost against mutual information about the latent regime and shared factor, with a closed-form shared-factor term and a lightweight Monte Carlo mode-identification term (Section~\ref{sec:exploration}; Appendix~\ref{app:info_gain}).
    \item We empirically evaluate on a synthetic correlated-regime simulator, on Temperature in Melbourne \citep{daily_min_temp_melbourne_kaggle}, and on Jena Climate (plus FRED DGS10 in the appendix), comparing against contextual bandit baselines, non-stationary bandit baselines, and ablations of the shared factor, regime count, transition model, and EM estimation (Section~\ref{sec:experiments}; Appendix~\ref{sec:experiments-details}).
\end{itemize}



%\paragraph{Contributions.}
%Our main contributions are:
%\begin{itemize}
%    \item A sequential expert-routing formulation for non-stationary time series with partial feedback and a time-varying expert pool (Section~\ref{sec:problem_formulation}).
%    \item A factorized switching state-space model for expert residuals with context-dependent regime switching and a shared latent factor enabling cross-expert information transfer (Section~\ref{sec:generative_model}).
%    \item A scalable IMM-style filtering recursion with dynamic registry management that updates only the queried expert jointly with the shared factor, while keeping per-expert marginals (Sections~\ref{sec:generative_model} and~\ref{sec:registry}).
%    \item An information-directed routing rule based on mutual information about $(z_t,\mathbf{g}_t)$, together with a predictive scheduling layer that outputs on-call active sets under availability constraints (Sections~\ref{sec:exploration} and~\ref{sec:scheduling}).
%\end{itemize}
